\runninghead{One to His Farm, Another to His Merchandise}
\properlane{farm-and-merchandise}{One to His Farm, Another to His Merchandise: Why the Invited Would Not Come}
\label{sec:farm-and-merchandise}

The parable's strangest fact is the refusal itself, and Chrysostom asks the
question aloud: ``For who would not choose to come to a marriage, and that a King's
marriage, and of a King making a marriage for a Son?'' (\work{Homilies on
Matthew} 69.1). The first refusal in Matthew is not hostility. The invited
``neglected and went their ways, one to his farm and another to his
merchandise'' (Mt 22:5); only ``the rest'' seize and kill the servants (22:6).
This interpretation asks what kept the first ones away. The witnesses who
expound the parable answer each in his own terms. Gregory reads the farm and
the merchandise as earthly labour beyond measure and the greed of gain, and
Hilary as worldly ambition and the love of money; Chrysostom and Aquinas
judge that no temporal cause, however reasonable it seems, excuses a refusal
of God; Jerome weighs the busy against the violent. Augustine and Bellarmine
do not comment on the parable. They expound the psalm verse of the first
Communion antiphon, on the rich who went hungry while those who seek the Lord
lack nothing. The interpretation as a whole is the editor's: read along this
line, the Mass sets beside the preoccupied the apostle who has learned to be
full and to be hungry, and the rich who went hungry beside those who seek the
Lord. No witness joins the parable's excuses to Philippians or to the psalm
verse. Both forms of the Gospel contain the verses this interpretation rests
on.

\subsection{The farm and the merchandise}

Matthew gives two excuses, the farm and the merchandise. Luke's parable of the
great supper has three, a farm, five yoke of oxen and a wife (Lk 14:18--20).
John Chrysostom, preaching on Matthew, lists ``yokes of oxen, and pieces of
ground, and wives'', drawing on Luke; the oxen and the wife are not in this
Gospel. What he draws from the excuses holds for Matthew's two: ``And yet the
excuses seem to be reasonable; but hence we learn, though the things which
hinder us be necessary, to set the things spiritual at a higher price than
all.'' The invited did not stay away because they were too busy: ``For not for
press of business, but from `making light of it,' they did not come''; and
``when spiritual things call us, there is no press of business that has the
power of necessity'' (69.1). Thomas Aquinas makes the same judgment in a
sentence: \latin{Videbantur habere iustam causam exterius; sed Dominus non
recipit: quia nulla temporalia debent detinere de veniendo ad Deum}, outwardly
they seemed to have a just cause, but the Lord does not accept it, because no
temporal things ought to keep anyone from coming to God (\work{Super Matthaeum}
c.~22, Venice, p.~281).

Gregory the Great reads the two excuses as two habits of life. \latin{In
villam quippe ire est labori terreno immoderate incumbere, in negotiationem
vero ire est actionum saecularium lucris inhiare}: to go to the farm is to lie
on earthly labour beyond measure, and to go to the merchandise is to gape
after the profits of worldly business. What they refuse is not a meal but a
mystery. The man intent on earthly labour, or given over to the business of
this world, \latin{mysterium incarnationis dominicae pensare et secundum illud
vivere dissimulat}, pretends not to see the need to weigh the mystery of the
Lord's Incarnation and to live by it, and so, as if going off to farm or trade,
he refuses to come to the king's wedding (\work{Homiliae in Evangelia} 38.5).
Hilary of Poitiers reads the excuses as ambition and avarice: some were held
\latin{ambitione saeculi tamquam agro}, by worldly ambition as by a field, and
more \latin{ob pecuniae cupiditatem negotiatione}, by trading for love of money
(\work{Commentarius in Matthaeum} 22.5, PL 9, col.~1043). Aquinas reports
Hilary's reading, the farm as the appetite for human glory and the trade as
\latin{appetitus avaritiae}, the appetite of avarice (p.~281).

Jerome sets the busy beside the violent and weighs them: \latin{Minoris enim
criminis sunt qui occupati aliis rebus venire noluerunt, his qui, contempto
invitantis affectu, verterunt humanitatem in crudelitatem}, those who, taken up
with other things, would not come are guilty of a lesser crime than those who,
despising the affection of the one inviting, turned kindness into cruelty
(\work{Commentarii in Matthaeum} III, PL 26, col.~159). The lesser crime is
still a refusal, and the busy are among those who ``were not worthy'' (22:8).

\subsection{Those whom no prosperity accompanies}

The king's second command sends the servants \latin{ad exitus viarum}, to the
places where the roads go out. Gregory reads the roads as men's actions and
their outlets as the failure of actions, \latin{quia illi plerumque facile ad
Deum veniunt, quos in terrenis actibus prospera nulla comitantur}, because
those very often come easily to God whom no prosperity accompanies in their
earthly doings (38.6). Jerome reads the same words of the gentiles:
\latin{Gentilium populus non erat in viis, sed in exitibus viarum}, the people
of the gentiles was not on the roads but at the outlets of the roads
(col.~160). The two readings answer different questions, Jerome's who the new
guests were and Gregory's why they came; Gregory's reading closes the circle of
this interpretation, for those who came easily were those whom the farm and
the merchandise did not hold.

\subsection{To be full and to be hungry}

The second reading gives the other side of the refusal: a man who has learned
not to be held by either plenty or want. ``I know both how to be brought low,
and I know how to abound \ldots\ both to be full and to be hungry: both to
abound and to suffer need. I can do all things in him who strengtheneth me''
(Phil 4:12--13). John Chrysostom hears in ``I have learned'' a discipline:
``this is an object of discipline, and exercise, and care, for it is not easy of
attainment, but very difficult, and a new thing.'' Plenty, he says, is as
dangerous as want: ``There is great need of virtue, not less than in the other
case. For as want inclines us to do many evil things, so too doth plenty'';
``Many know not how to be full, as for example, the Israelites, `ate, and
kicked'\,'' (Deut 32:15). And he will not let Paul boast: ``The success is not
mine own, but His who has given me strength'' (\work{Homilies on Philippians}
15).

Thomas Aquinas makes the same verses the mark of the wise man.
\latin{Nihil demonstrat ita mentem sapientis perfecti, sicut quod sciat uti
quolibet statu}: nothing shows the mind of a perfect wise man so well as that he
knows how to use any state; \latin{ille perfectus est qui scit uti quolibet
statu, ut si sit in magno non elevetur, et si in minimo non deiiciatur}, he is
perfect who knows how to use any state, so that in a great one he is not lifted
up and in the least he is not cast down. Both states carry a danger: \latin{in
utroque est periculum, quia ex abundantia erigitur animus contra Deum, ex
paupertate deiicitur \ldots\ Sed Apostolus \ldots\ virtute in utroque scit uti},
in both there is danger, because from abundance the mind is raised against God
and from poverty it is cast down; but the Apostle knows how to use both with
virtue. And ``I can do all things'' means \latin{Non possem hos insultus
sustinere, nisi manu Dei me confortante}, I could not bear these assaults unless
the hand of God were strengthening me (\work{Super Philippenses} c.~4, lect.~2,
Dessain, pp.~401--403).

Heard beside the parable, Paul is the guest who would have come. He is not
without goods, for he has just received the Philippians' renewed care (4:10,
14); he is free of being held by them. The invited went off to farm and trade, and Paul, who knows
abundance and want alike, says he can do all things in the one who strengthens
him. The \work{Ordo}'s title for the reading,
\latin{Omnia possum in eo, qui me confortat}, singles out the verse on which it
turns, within the reading's own course.

The reading ends with God's supplying of the Philippians' want, and the two
witnesses read that verse from different texts. In the Vulgate and the
Douay--Rheims it is a wish, ``may my God supply all your want'', and Aquinas
expounds the Vulgate's \latin{impleat omne desiderium}: first as God's return
to the givers who had filled Paul's need, \latin{Impleat, etc., quia implestis
meum}, may he fill yours, because you have filled mine; then of the fulfilment
of all desire in glory, \latin{quia ibi implebitur totum desiderium}, with the
psalm verse \latin{Satiabor cum apparuerit gloria tua} (p.~403). Chrysostom's
text is a blessing, and he notes that Paul ``invokes blessings upon them, as poor
men do''. The \work{Nova Vulgata}, which the \work{Ordo} follows, has a
promise, \latin{implebit omne desiderium}, with the same \latin{desiderium}:
the two Latin texts differ only in the mood of the verb. Aquinas's reading of
desire filled in glory rests on that word, desire, which the Douay--Rheims
renders ``want''; it is used here as his.

\subsection{The rich who went hungry}

Where the first Communion antiphon is chosen, the Mass ends this line with the
psalm verse on the rich who hungered. Augustine meets the objection that the
verse seems false, for ``many rich men that are wicked die in their riches'',
and the man who sees it says, ``the Scripture has deceived me''; such a man
``seeketh mortal food on the earth, and seeketh not a true reward in heaven''.
The answer is the spiritual sense. ``When thou art filled with spiritual riches,
canst thou be poor? And was he therefore rich, because he had a bed of ivory;
and art thou poor who hast the chamber of thy heart filled with such jewelry of
virtues, justice, truth, charity, faith, endurance?'' The rich who seem full
are the hungry ones: ``Those rich then lack, they lack, and what is heavier,
they lack bread'', the bread of the one who said, ``I am the Living Bread which
came down from Heaven'' (\work{Enarrationes in Psalmos} 33, sections 13--14 of
the English text). Robert Bellarmine gives the verse a literal sense, that the
rich ``began to hunger and to need, because riches are fallacious and
uncertain'', and then ``a higher meaning'': ``those who are attached to the
temporalities of this world always hunger and need, for they are always
covetous and desirous of having more; but `they that seek the Lord,' as they
seek a thing of infinite value, a thing greater than their desires \ldots\ as
they cling to the Supreme Good, they possess all that is good''
(\work{Commentary on the Book of Psalms}, Ps~33:10--11).

The Hebrew of the verse has young lions where the Latin has the rich, so this
reading rests on the Latin and Greek text that the Missal cites. Within that
text, the antiphon answers the parable's refusers exactly: they went off to
their goods and are the rich who go hungry; the seekers of the Lord lack
nothing. The responsorial psalm has already said the same in its first line,
``The Lord ruleth me: and I shall want nothing'', which Aquinas, as quoted with
the psalm above, reads of the future, when nothing will be lacking because we
shall have God.

\subsection{The Missal's texts in this reading}

The Entrance Antiphon's question, who shall stand, meets in Augustine's
exposition, quoted with the antiphon above, a description of the preoccupied:
those who prosper in iniquity are the deeper in the depths the happier they
seem, for ``a deceitful happiness is itself a greater unhappiness''. The Collect
asks that grace keep the assembly \latin{intentos}, intent, on good works; the
word names the attention that the invited gave to farm and trade.

The first reading supplies what the refusers did not do. ``Lo, this is our
God, we have waited for him'' (Is 25:9), and Jerome glosses the waiting:
\latin{exspectavimus eum, hoc est, verbis eius credidimus}, we waited for him,
that is, we believed his words (\work{Commentarii in Isaiam} VIII, PL 24,
cols.~300--301). The invited of the parable had been told twice; the people of
Isaiah's mountain believed the word and waited.
The acclamation asks for eyes enlightened to know the hope of our calling, and
Aquinas, as quoted with the acclamation above, reads Paul's prayer against
those whose eyes are enlightened only to know temporal things. The Prayer
over the Offerings offers the faithful's prayers with the sacrificial gifts:
heard in this reading, goods given rather than held. Chrysostom's word on the
Philippians' gift, that what counted was ``not the money, but the purpose'',
belongs to v.~18, which is not read, and stands here only as context.

Where the second Communion antiphon is chosen, the promise of seeing him as he
is becomes, in Augustine's exposition of the verse, a lesson in desire. ``The
whole life of a good Christian is an holy desire''; God, ``by deferring our
hope, stretches our desire; by the desiring, stretches the mind; by
stretching, makes it more capacious. Let us desire therefore, my brethren, for
we shall be filled.'' But the vessel must first be emptied: ``holy longing
exercises us just so much as we prune off our longings from the love of the
world. \ldots\ Suppose that God would fill thee with honey: if thou art full of
vinegar, where wilt thou put the honey?'' (\work{Homilies on the First Epistle
of John} 4.6). The guests who went off to their farms were full. The Prayer
after Communion asks to be made partakers of the divine nature, and the verse
of 2~Peter whose words it uses continues, ``flying the corruption of that
concupiscence which is in the world'' (2~Pet 1:4), though the prayer does not
quote that clause.

\subsection{Where this interpretation strains}

The refusers in the parable are not only busy. ``The rest'' seize the
servants, treat them shamefully and kill them (22:6), and Jerome's distinction
between the lesser crime of the occupied and the greater crime of the cruel
shows that a reading of the parable as a story of distraction misses half of
it; the first interpretation keeps that half. Matthew's excuses are two, and
Chrysostom's oxen and wife belong to Luke. The second reading follows its own
course through Philippians and comes to its end this Sunday for that reason;
no witness joins it to the parable or sets Paul beside the invited. The first Communion antiphon's rich are the reading
of the Latin and Greek text, not of the Hebrew. And the verse on God's
supplying of want is a wish in the Vulgate and a promise in the \work{Nova
Vulgata}, though both read \latin{desiderium}; what Aquinas draws from that
word is drawn from the Latin, and an English that speaks of want or need
does not carry it.

\Needspace{26\baselineskip}
\subsection{The four senses of this interpretation}

\begin{fourSenses}
\item[Literal.] The invited make light of the king's wedding and go off, one to
his farm and another to his merchandise. Paul writes that he has learned to be
full and to be hungry, to abound and to suffer need, and that he can do all
things in him who strengthens him.

\item[Allegorical.] The feast refused is the wedding in which the Father joined
the holy Church to the Son through the mystery of the Incarnation (Gregory),
the mystery the preoccupied will not weigh or live by (Gregory). Refused, the
call passes to those at the outlets of the roads, whom no earthly prosperity
accompanies (Gregory), the people of the gentiles (Jerome).

\item[Moral.] The farm and the merchandise are earthly labour without measure
and the hunger for gain (Gregory), worldly ambition and the love of money
(Hilary). Set spiritual things above even necessary ones (Chrysostom); use
plenty and want alike with virtue, by Christ's strength (Aquinas). Where the
first Communion antiphon is chosen: seek the Lord, who fills what riches leave
hungry (Augustine, Bellarmine).

\item[Anagogical.] God fills all desire in glory: ``I shall be satisfied when
thy glory shall appear'' (Aquinas, on the \latin{desiderium} of
Philippians 4:19). Where the second Communion antiphon is chosen: the desire stretched
by waiting will be filled by the vision (Augustine).
\end{fourSenses}
